\chapter{Introdução}

A construção de ambientes virtuais para robótica não termina quando uma cena reconstruída produz imagens semelhantes às observações originais. Uma reconstrução pode representar corretamente cores, contornos e iluminação nos pontos de vista avaliados e, ainda assim, posicionar uma superfície alguns centímetros distante de sua posição de referência. Essa diferença tem efeito direto quando o simulador calcula o contato de uma roda, de uma caixa ou de um robô com o ambiente.

A organização da geometria também é necessária. Se uma mesa, uma cadeira e o piso permanecem unidos em uma única malha, o simulador consegue exibir o conjunto, mas não consegue deslocar somente a cadeira sem processamento adicional. Uma superfície que se parece com um piso nas imagens pode conter ondulações ou vértices isolados que alteram as colisões. Portanto, produzir novas imagens, reconstruir a superfície e organizar a cena em objetos utilizáveis pelo Gazebo são problemas relacionados, porém distintos.

Os campos de radiância neurais, introduzidos por \citet{mildenhall2020nerf}, aprendem essa representação a partir de imagens e das posições e orientações das câmeras. Durante o treinamento, uma rede relaciona pontos do espaço às cores e à ocupação observadas nas imagens. Depois do treinamento, a representação pode ser consultada a partir de uma nova câmera para gerar uma imagem do ambiente. Esse processo reduz a modelagem manual da aparência, mas sua saída continua sendo uma função aprendida, e não um conjunto de modelos separados com geometrias visuais e de colisão.

O DM-NeRF \citep{wang2023dmnerf} acrescenta a essa representação a identificação das instâncias. Uma instância corresponde a um objeto individual: duas cadeiras, por exemplo, recebem identificadores diferentes mesmo que pertençam à mesma categoria. No protocolo DM-SR utilizado nesta dissertação, o treinamento recebe imagens RGB, poses de câmera e mapas bidimensionais de instância. Cada pixel desses mapas contém um rótulo que informa qual objeto aparece naquela posição da imagem. As máscaras são fornecidas ao treinamento; portanto, a avaliação não mede descoberta de objetos sem supervisão.

O DM-NeRF representa a cena e permite consultar, em cada posição do espaço, tanto informações de aparência quanto a identidade de instância. Para chegar ao Gazebo, ainda é preciso converter essa representação em uma malha triangular, associar os identificadores aos seus vértices e separar os triângulos em arquivos de objetos independentes. Também é preciso identificar uma superfície de apoio, tratar sua geometria quando ela satisfaz os critérios definidos e criar uma descrição de cena compatível com o simulador.

A dissertação examina essa passagem da representação aprendida pelo DM-NeRF para uma cena explícita, organizada em objetos independentes e utilizável no Gazebo Classic. A cadeia estudada compreende: (1) a reprodução do DM-NeRF na cena \textit{study} do DM-SR; (2) a extração de uma malha; (3) a associação explícita dos identificadores de instância aos vértices; (4) a separação dos objetos; (5) a identificação geométrica de um candidato a piso; (6) a regularização condicionada dessa superfície; e (7) a exportação para Gazebo Classic. O SDFormat é empregado para descrever os modelos e o mundo carregado pelo simulador.

\section{Problema de pesquisa}

Simuladores como o Gazebo precisam que os elementos da cena sejam definidos individualmente, com geometria, posição e propriedades de colisão. Uma mesa, por exemplo, precisa corresponder a um modelo identificável para que sua pose possa ser consultada ou alterada sem mover o restante do ambiente. Quando toda a reconstrução forma uma única malha, essa operação não é possível diretamente.

Há ainda uma diferença entre aparência e geometria. Uma pequena alteração de profundidade pode quase não ser percebida em uma imagem renderizada, mas pode mudar o ponto no qual dois corpos colidem. Da mesma forma, uma superfície de apoio irregular pode sustentar um objeto em uma posição diferente da esperada ou permitir que ele atravesse uma região da malha. Além disso, massa, inércia, atrito e articulações não são obtidos de maneira geral apenas a partir das imagens usadas neste estudo.

A questão que orienta a dissertação é:

\begin{quote}
\textbf{Em que medida uma reconstrução neural com decomposição de instâncias, combinada a pós-processamento geométrico, permite gerar ambientes domésticos estruturados para simulação robótica?}
\end{quote}

Nesse enunciado, reconstrução neural designa a representação tridimensional aprendida pelo DM-NeRF a partir das imagens e poses. A resposta é construída por avaliações separadas: as imagens sintetizadas são comparadas às imagens de referência; as instâncias são comparadas por suas máscaras; a malha é comparada à geometria de referência; os objetos separados são verificados por sua topologia; e a cena exportada é carregada e testada no Gazebo.

\section{Objetivos}

O objetivo geral é avaliar a capacidade de uma cadeia baseada em DM-NeRF e pós-processamento geométrico de transformar uma representação aprendida da cena, com identificadores de instância, em um ambiente doméstico estruturado para o Gazebo Classic.

Os objetivos específicos são:
\begin{enumerate}[label=\arabic*)]
  \item reproduzir o DM-NeRF na cena \textit{study} do DM-SR e avaliar a síntese de vistas e a decomposição de instâncias pelo protocolo oficial;
  \item avaliar a qualidade geométrica da malha extraída, distinguindo-a da qualidade das imagens sintetizadas;
  \item estruturar a cena em objetos tridimensionais independentes e selecionar e regularizar geometricamente sua superfície de apoio; e
  \item avaliar a integração da cena estruturada no Gazebo Classic segundo critérios funcionais definidos previamente.
\end{enumerate}

O registro do custo computacional e da origem dos artefatos acompanha transversalmente essas etapas, de modo que as medições possam ser relacionadas às entradas e aos programas que as produziram.

\section{Delimitação e contribuição}

O escopo experimental adota o DM-NeRF como método principal, o DM-SR como conjunto de dados e a cena \textit{study} como estudo de caso. A avaliação geométrica inclui o COLMAP, que reconstrói a superfície comparando imagens obtidas de diferentes posições. Ele funciona como referência externa para a geometria tridimensional; a identificação de instâncias e a integração com o Gazebo são avaliadas somente na cadeia baseada em DM-NeRF.

A contribuição consiste na formulação e na análise experimental do percurso entre a saída do DM-NeRF e uma cena estruturada para simulação. O trabalho associa os identificadores de instância à topologia da malha, trata separadamente as faces situadas entre rótulos diferentes, seleciona uma superfície de apoio por critérios geométricos e a regulariza somente quando esses critérios são satisfeitos. Cada resultado é medido no ponto correspondente da cadeia: aparência e máscaras nas vistas, distâncias na geometria, preservação de faces na separação e funcionamento após a exportação.

Neste experimento, o DM-NeRF fornece identificadores que permitem distinguir os objetos individuais da cena. Esses identificadores são usados para separar os objetos e organizar a cena para simulação. A avaliação verifica essa separação; identificar a categoria de cada objeto, como ``cadeira'' ou ``mesa'', fica fora do escopo.

O processo recebe como entrada imagens, posições e orientações das câmeras e mapas que identificam os objetos em cada imagem. Com esses dados, o DM-NeRF é treinado para representar a cena e suas instâncias. Os programas consultam o modelo treinado para extrair uma malha e a separam em objetos individuais. Em seguida, procuram uma superfície que possa corresponder ao piso e, quando ela atende aos critérios definidos, ajustam seus vértices a um plano. Por fim, geram os arquivos que carregam a cena no Gazebo. É essa sequência de processamento que o trabalho chama de automatizada: depois que as entradas são fornecidas, as etapas são executadas sem editar manualmente a malha de cada objeto.
